\paragraph{DDPO and the reference reward}
Denoising Diffusion Policy Optimization (DDPO) \citep{black2024ddpo} views a denoising trajectory as an episode. The state contains the noisy latent, timestep and prompt. The action is the next latent. A frozen scorer evaluates the decoded image at the end.

\paragraph{How a final reward trains earlier steps.}
Retain the previous section's latent dimension $d\in\mathbb N_{>0}$, conditioning $c\in\mathcal C$ and parameter vector $\theta\in\mathbb R^{d_\theta}$, with $d_\theta\in\mathbb N_{>0}$. An inference schedule has $J\in\mathbb N_{>0}$ transitions at integer indices $t_0=T>t_1>\cdots>t_J=0$, where $T\in\mathbb N_{>0}$ is the training noise horizon. The trajectory is $\xi=(x_{t_0},\ldots,x_{t_J})\in(\mathbb R^d)^{J+1}$. Let $\mathcal S\subseteq\{1,\ldots,J\}$ index its positive-variance transitions. Their scalar log densities are
$\ell_{\theta,j}=\log p_\theta(x_{t_j}\mid x_{t_{j-1}},c)\in\mathbb R$.
Let $R(\xi,c)$ be the terminal reward, $b(c)$ a prompt-dependent baseline and $P_\theta(\xi\mid c)$ the trajectory density in the all-stochastic case. With fixed initial noise law, fixed reward as a function of the sampled trajectory, and differentiation under the integral,
\begin{align}
 \nabla_\theta\E_{P_\theta}[R(\xi,c)]
   &=\int R(\xi,c)\,\nabla_\theta P_\theta(\xi\mid c)\,d\xi
       \label{eq:ddpo-differentiate-expectation}\\
   &=\E_{P_\theta}\!\left[R(\xi,c)
       \nabla_\theta\log P_\theta(\xi\mid c)\right]
       \label{eq:ddpo-score-identity}\\
   &=\E_{P_\theta}\!\left[(R(\xi,c)-b(c))
           \sum_{j\in\mathcal S}\nabla_\theta\ell_{\theta,j}\right].
       \label{eq:ddpo-policy-gradient}
\end{align}
The trajectory density factors into the initial noise density and the transition densities. The baseline term vanishes because the integral of the trajectory density is one for every $\theta$. A positive $R-b$ increases the likelihood of the sampled transitions. The local trainer omits deterministic terminal steps and holds saved rewards fixed, so its update is a surrogate if a skipped terminal map depends on $\theta$.

\paragraph{Reusing trajectories and measuring drift.}
For a saved stochastic transition, the current-to-rollout density ratio is
$u_j=\exp(\ell_{\theta,j}-\ell_{{\rm old},j})\in(0,\infty)$, with $\ell_{{\rm old},j}\in\mathbb R$ recorded at generation. It evaluates the same saved successor under both models. DDPO uses such ratios with clipping, analogously to GRPO. The local implementation averages log densities over coordinates and therefore uses $\exp[(\ell_{\theta,j}-\ell_{{\rm old},j})/d]$. This changes the scale on which clipping acts.

At one state, write $p_\theta$ and $p_{\rm ref}$ for current and frozen-reference Gaussian transition laws on $\mathbb R^d$. With mean vectors $\mu_\theta,\mu_{\rm ref}\in\mathbb R^d$ and equal scalar variance $\sigma^2\in(0,\infty)$, their scalar conditional divergence is
\begin{equation}
\label{eq:gaussian-transition-kl}
 \KL(p_\theta\Vert p_{\rm ref})
 =\frac{\|\mu_\theta-\mu_{\rm ref}\|_2^2}{2\sigma^2}\in[0,\infty).
\end{equation}
The project averages this divergence over coordinates and stochastic steps to obtain $\bar K\in[0,\infty)$, then uses reference reward $1/(1+\bar K/s_0)\in(0,1]$ for scale $s_0\in(0,\infty)$ fitted on training probes. Similar transition means earn a higher reward. Because a weighted reward can trade reference similarity for evasion, held-out prompt alignment is also measured directly.
